%%%% ABSTRACT
%%
%% Versão do resumo para idioma de divulgação internacional.

\begin{abstractutfpr}%% Ambiente abstractutfpr
In software development, it is common to use metrics to assess code quality and the time spent on its production. In the professional environment, this analysis occurs mainly through pull requests, while in the academic context, the assessment is usually done through occasional submissions (such as emails or learning management systems - LMS). Although LMSs record activity and event logs, there is a gap in the structured collection of
data on the programming learning process, especially in environments such as computer labs. This work proposes the creation of a dataset with logs of the code development process, capturing the learning trail in real time.
Unlike conventional approaches, which depend on Learning Management Systems Unlike conventional approaches, which rely on Learning Management Systems, data collection will be performed through a plugin integrated into a
widely used code editor, designed for both programming and documentation . It is expected to provide a database on programming teaching and learning that can assist in identifying student difficulties, personalizing teaching,
and improving pedagogical practices.

\end{abstractutfpr}
